\documentclass[12pt,a4paper]{article}

% ============================================================
% PAQUETES
% ============================================================
\usepackage[utf8]{inputenc}
\usepackage[T1]{fontenc}
\usepackage[spanish,es-nodecimaldot]{babel}
\usepackage{amsmath,amssymb,amsthm}
\usepackage{graphicx}
\usepackage{booktabs}
\usepackage{multirow}
\usepackage{array}
\usepackage{longtable}
\usepackage{tabularx}
\usepackage{float}
\usepackage{hyperref}
\usepackage{cleveref}
\usepackage{enumitem}
\usepackage{siunitx}
\usepackage{xcolor}
\usepackage{listings}
\usepackage{algorithm}
\usepackage{algpseudocode}
\usepackage{natbib}
\usepackage{fancyhdr}
\usepackage{geometry}
\usepackage{setspace}
\usepackage{titlesec}
\usepackage{tikz}
\usepackage{pgfplots}
\pgfplotsset{compat=1.18}
\usetikzlibrary{shapes,arrows,positioning,calc}

% ============================================================
% CONFIGURACIÓN DE PÁGINA
% ============================================================
\geometry{left=2.5cm,right=2.5cm,top=2.5cm,bottom=2.5cm}
\setstretch{1.15}

% ============================================================
% CONFIGURACIÓN DE ENCABEZADOS
% ============================================================
\pagestyle{fancy}
\fancyhf{}
\fancyhead[L]{\small Informe de Investigación -- BIP 2026}
\fancyhead[R]{\small\nouppercase{\leftmark}}
\fancyfoot[C]{\thepage}
\renewcommand{\headrulewidth}{0.4pt}

% ============================================================
% CONFIGURACIÓN DE TÍTULOS
% ============================================================
\titleformat{\section}{\Large\bfseries}{\thesection.}{0.5em}{}
\titleformat{\subsection}{\large\bfseries}{\thesubsection.}{0.5em}{}
\titleformat{\subsubsection}{\normalsize\bfseries}{\thesubsubsection.}{0.5em}{}

% ============================================================
% COLORES Y LISTADOS
% ============================================================
\definecolor{codegreen}{rgb}{0.25,0.49,0.48}
\definecolor{codegray}{rgb}{0.5,0.5,0.5}
\definecolor{codepurple}{rgb}{0.58,0,0.82}
\definecolor{backcolour}{rgb}{0.97,0.97,0.95}
\lstdefinestyle{mystyle}{
    backgroundcolor=\color{backcolour},
    commentstyle=\color{codegreen},
    keywordstyle=\color{magenta},
    numberstyle=\tiny\color{codegray},
    stringstyle=\color{codepurple},
    basicstyle=\ttfamily\footnotesize,
    breakatwhitespace=false,
    breaklines=true,
    captionpos=b,
    keepspaces=true,
    numbers=left,
    numbersep=5pt,
    showspaces=false,
    showstringspaces=false,
    showtabs=false,
    tabsize=2
}
\lstset{style=mystyle}

% ============================================================
% NUEVOS COMANDOS
% ============================================================
\newcommand{\KL}{\mathrm{KL}}
\newcommand{\JSD}{\mathrm{JSD}}
\newcommand{\AUROC}{\mathrm{AUROC}}
\newcommand{\AUPRC}{\mathrm{AUPRC}}
\newcommand{\ECE}{\mathrm{ECE}}
\newcommand{\VLM}{\text{VLM}}
\newcommand{\UQ}{\text{UQ}}
\newcommand{\XAI}{\text{XAI}}
\newcommand{\MedVLM}{\text{MedVLM}}
\DeclareMathOperator*{\argmax}{arg\,max}
\DeclareMathOperator*{\argmin}{arg\,min}

% ============================================================
% INFORMACIÓN DEL DOCUMENTO
% ============================================================
\title{\textbf{Cross-Modal Representation Disagreement as a Lightweight Uncertainty Signal for Glaucoma Detection in Medical Vision-Language Models}}
\author{
    \textbf{Miguel Guillermo Abreu Cárdenas}\\[0.3em]
    Doctorado en Inteligencia Artificial Aplicada a Oftalmología\\
    Tutor: Dr.~Saúl Calderón Ramírez\\[0.5em]
    \textit{Instituto Tecnológico de Costa Rica}
}
\date{Julio 2026}

% ============================================================
% INICIO DEL DOCUMENTO
% ============================================================
\begin{document}

% ---------- PORTADA ----------
\begin{titlepage}
    \centering
    \vspace*{1cm}
    
    {\Large\textsc{Instituto Tecnológico de Costa Rica}}\\[0.3cm]
    {\large Programa de Doctorado en Inteligencia Artificial}\\[0.3cm]
    {\large Área de Investigación: IA Aplicada a Oftalmología}
    
    \vspace{2cm}
    
    \rule{\linewidth}{0.5mm}\\[0.6cm]
    {\Huge\bfseries Cross-Modal Representation Disagreement as a Lightweight Uncertainty Signal for Glaucoma Detection in Medical Vision-Language Models}\\[0.3cm]
    \rule{\linewidth}{0.5mm}
    
    \vspace{1.5cm}
    
    {\Large\textbf{Informe de Investigación Académica}}\\[0.5cm]
    {\large Propuesta de Paper para BIP 2026}
    
    \vspace{2cm}
    
    \begin{tabular}{rl}
        \textbf{Autor:} & Miguel Guillermo Abreu Cárdenas \\[0.3em]
        \textbf{Tutor:} & Dr.~Saúl Calderón Ramírez \\[0.3em]
        \textbf{Fecha:} & Julio 2026 \\[0.3em]
        \textbf{Venue Objetivo:} & 8th IEEE International Conference on BioInspired Processing (BIP 2026) \\
    \end{tabular}
    
    \vfill
    
    {\small Costa Rica, 11--13 de Noviembre de 2026}
\end{titlepage}

% ---------- ÍNDICE ----------
\newpage
\tableofcontents
\newpage

% ============================================================
\section{Resumen Ejecutivo}
% ============================================================

Este informe consolida la investigación desarrollada para la propuesta del paper \textit{Cross-Modal Representation Disagreement as a Lightweight Uncertainty Signal for Glaucoma Detection in Medical Vision-Language Models}, destinado a la 8th IEEE International Conference on BioInspired Processing (BIP 2026).

\subsection{La contribución}

Este paper propone que la \textbf{divergencia KL entre las representaciones visuales y textuales} de un Medical Vision-Language Model (MedGemma 4B), operando en \textbf{un solo forward pass} y sin entrenamiento adicional, constituye una señal de incertidumbre \textit{lightweight}, interpretable y clínicamente accionable para la detección de glaucoma en imágenes de fondo de ojo.

A diferencia de métodos como MC-Dropout (multi-pass, costoso) o entropía predictiva (solo textual, no detecta alucinaciones visuales), nuestra señal captura el \textbf{desacuerdo cross-modal} entre lo que el modelo ``ve'' y lo que el modelo ``dice''.
Evaluamos esta hipótesis en el dataset MM-ODIR-129 (129 imágenes de fondo de ojo, 60 normales y 69 patológicas con glaucoma), comparando contra múltiples baselines de \UQ\ y reportando \AUROC, \AUPRC, Brier score y \ECE.

\subsection{Relevancia para la tesis doctoral}

Los cuatro pilares de la tesis doctoral se articulan naturalmente alrededor de este método:

\begin{enumerate}
    \item \textbf{UQ (Pilar 1):} La KL cross-modal es una nueva familia de métodos \textit{single-pass} + \textit{feature-based} + \textit{cross-modal} que llena un vacío en el espacio metodológico.
    
    \item \textbf{XAI (Pilar 2):} La dirección de la KL es inherentemente explicativa: diagnostica \textit{qué tipo} de incertidumbre existe (visión vs.~texto).
    
    \item \textbf{Few-Shot (Pilar 3):} La KL puede filtrar pseudo-labels y seleccionar support sets de alta calidad en adaptación few-shot.
    
    \item \textbf{Segmentación (Pilar 4):} La KL global puede descomponerse en mapas de riesgo espacial para segmentación de disco/copa óptica.
\end{enumerate}

\subsection{Hallazgos clave}

El trabajo consolidado en este informe abarca diez dimensiones de investigación profunda, cross-verificación de hallazgos y nueve insights cross-dimensionales.
La metodología propuesta combina tres propiedades que ningún método existente une simultáneamente: \textbf{single-pass}, \textbf{basada en features internas} y \textbf{cross-modal}.

% ============================================================
\section{Introducción}
% ============================================================

\subsection{Contexto general}

El glaucoma constituye la principal causa irreversible de ceguera a nivel mundial, afectando a más de 80 millones de personas.
Su detección temprana mediante el análisis de imágenes de fondo de ojo (fundus photography) representa una de las aplicaciones más prometedoras de la inteligencia artificial en salud.
Sin embargo, los modelos de deep learning tradicionales, aunque alcanzan alta accuracy en benchmarks controlados, frecuentemente exhiben predicciones de \textbf{alta confianza en casos erróneos} (\textit{overconfident errors}), lo que los hace inseguros para despliegue clínico.

\subsection{El problema de la incertidumbre en modelos médicos}

La estimación de incertidumbre (\UQ, \textit{Uncertainty Quantification}) en modelos de visión-lenguaje médicos (\MedVLM s) enfrenta desafíos únicos:

\begin{itemize}[leftmargin=2em]
    \item \textbf{Modality bias:} Los \VLM s clínicos frecuentemente favorecen priores lingüísticos sobre evidencia visual, generando predicciones plausibles pero alucinadas \citep{modalitybias2025}.
    
    \item \textbf{Costo computacional:} Métodos como MC-Dropout requieren 25--100 pasadas por imagen, y los deep ensembles requieren múltiples modelos en memoria \citep{gal2016dropout, lakshminarayanan2017ensembles}.
    
    \item \textbf{Limitaciones de baselines baratos:} La entropía predictiva y el Maximum Softmax Probability (MSP) solo miran la salida textual, sin detectar cuando el modelo \textit{alucina} ignorando la imagen.
\end{itemize}

\subsection{Motivación bio-inspirada}

El cerebro humano detecta conflictos entre modalidades sensoriales como señal de alerta.
El \textbf{efecto McGurk} demuestra que cuando la información visual (movimientos labiales) y auditiva (sonido) discrepan, el cerebro registra el conflicto \citep{mcgurk1976}.
Esta detección de conflictos cross-modal es un principio neurocientífico bien establecido que motiva nuestra aproximación computacional.

\subsection{Propuesta}

Proponemos medir la divergencia KL entre las representaciones internas visuales y textuales de MedGemma 4B en una \textbf{única pasada forward}.
Nuestra hipótesis central es:

\begin{quote}
\textbf{H1:} Cuando el modelo se equivoca al diagnosticar glaucoma, su desacuerdo interno $u(x) = \KL(p_{\text{vis}} \| p_{\text{text}})$ es significativamente mayor que cuando acierta.
\end{quote}

Adicionalmente, investigamos una hipótesis secundaria:

\begin{quote}
\textbf{H4:} En casos patológicos, existe correlación entre la severidad del daño (medida por cup-to-disc ratio) y la magnitud de la señal de incertidumbre.
\end{quote}

% ============================================================
\section{Estado del Arte y Análisis de la Literatura}
% ============================================================

\subsection{Vision-Language Models en oftalmología}

Los modelos fundacionales multimodales han revolucionado el análisis de imágenes médicas.
En oftalmología, destacan:

\begin{itemize}[leftmargin=2em]
    \item \textbf{EyeCLIP} \citep{shietal2025eyeclip}: Modelo visual-lenguaje multimodal específico para oftalmología computacional.
    
    \item \textbf{RetiZero} \citep{retizero2024}: Identificación de enfermedades comunes y raras del fondo de ojo usando modelos fundacionales VLM con conocimiento de más de 400 enfermedades.
    
    \item \textbf{VOLMO} \citep{volmo2026}: Modelos versátiles y abiertos para oftalmología.
    
    \item \textbf{MedGemma} \citep{sellergren2025medgemma}: Familia de modelos fundacionales médicos de Google, con variantes de 4B parámetros especializadas en visión-lenguaje médico.
\end{itemize}

MedGemma-4B alcanza un F1 de aproximadamente 63.7\% en glaucoma \citep{volmo2026}, lo que lo hace un clasificador base modesto pero ideal para evaluar \UQ: genera abundantes errores con los que validar la señal de incertidumbre.

\subsection{Métodos de estimación de incertidumbre}

\subsubsection{Métodos multi-pass (costosos)}

\begin{itemize}[leftmargin=2em]
    \item \textbf{MC-Dropout} \citep{gal2016dropout}: Múltiples inferencias con dropout activo.
    Requiere 25--100 pasadas.
    
    \item \textbf{Deep Ensembles} \citep{lakshminarayanan2017ensembles}: Múltiples modelos entrenados con diferentes inicializaciones.
    Requiere 5--10 modelos en memoria.
    
    \item \textbf{Semantic Entropy} \citep{kuhn2023semantic}: Múltiples generaciones muestreadas + modelo NLI externo para agrupar respuestas semánticamente equivalentes.
    Requiere 10+ generaciones.
\end{itemize}

\subsubsection{Métodos single-pass}

\begin{itemize}[leftmargin=2em]
    \item \textbf{Maximum Softmax Probability (MSP)} \citep{hendrycks2017baseline}: La máxima probabilidad de la distribución de salida.
    
    \item \textbf{Entropía predictiva}: La entropía de la distribución de salida sobre el vocabulario.
    
    \item \textbf{Energy Score} \citep{liu2020energy}: Basado en la función de energía de la distribución de salida.
    
    \item \textbf{Mahalanobis Distance} \citep{lee2018mahalanobis}: Distancia en el espacio de features respecto a la distribución de clases.
\end{itemize}

\subsubsection{Métodos basados en features internas}

\begin{itemize}[leftmargin=2em]
    \item \textbf{SAPLMA} \citep{azaria2023internal}: MLP sobre hidden states para predecir veracidad.
    Requiere entrenar un probe con datos etiquetados.
    
    \item \textbf{``Between the Layers Lies the Truth''} \citep{badash2026between}: Convierte activaciones post-MLP de cada capa en distribuciones vía softmax temperado y computa matriz de divergencias KL pairwise entre capas.
    Entrena un clasificador LightGBM sobre la ``signature map'' aplanada.
    AUPRC $\sim$0.82 bajo cuantización 4-bit.
\end{itemize}

\subsection{Métodos cross-modal de incertidumbre}

\begin{itemize}[leftmargin=2em]
    \item \textbf{VLM-UQBench} \citep{vlmuqbench2026}: Benchmark de 600 muestras evaluando 9 métodos de incertidumbre visual, textual y cross-modal.
    
    \item \textbf{UMPIRE} \citep{umpire2026}: Framework training-free usando incoherence-adjusted semantic volume con DPPs.
    Es multi-sample, no single-pass.
    
    \item \textbf{VIG-TUQ} \citep{vigtuq2026}: Visual-Grounded Token Uncertainty.
    Usa Jensen-Shannon entre predicciones con/sin imagen + attention weights.
    \textbf{Nota crítica:} su score JSD requiere un segundo forward \textit{sin} imagen; solo su score de atención es single-pass.
    
    \item \textbf{Expert-CFG} \citep{liang2025uncertainty}: Usa entropía predictiva para detectar respuestas no confiables en MedVLMs, luego aplica classifier-free guidance con anotaciones expertas.
    AUC $>$ 0.8 en VQA-RAD/SLAKE.
\end{itemize}

\subsection{El ``punto dulce'' metodológico}

\begin{figure}[H]
\centering
\begin{tikzpicture}[
    axis/.style={->,>=stealth,thick},
    node distance=1.2cm,
    box/.style={draw, rounded corners, minimum width=3cm, minimum height=0.8cm, align=center, font=\small}
]
% Ejes
\draw[axis] (0,0) -- (5,0) node[right] {\small Costo computacional};
\draw[axis] (0,0) -- (0,4) node[above] {\small Información cross-modal};
\draw[axis] (0,0) -- (-3,-2) node[below left] {\small Features internas};

% Métodos
\node[box, fill=red!15] at (4,0.5) {MSP/Entropy};
\node[box, fill=red!15] at (4,1.5) {Energy/Mahalanobis};
\node[box, fill=orange!20] at (1,3) {Semantic Entropy};
\node[box, fill=orange!20] at (2,3.5) {MC-Dropout};
\node[box, fill=green!20] at (-1.5,2) {SAPLMA};
\node[box, fill=green!30, very thick] at (-1,1) {\textbf{KL Cross-Modal}};
\node[box, fill=blue!15] at (-2,3) {Between the Layers};

% Leyenda
\node[draw, rounded corners, fill=green!10, font=\footnotesize, align=left] at (3,-1.5) {
    \textbf{Nuestra contribución}\\
    Single-pass + Feature-based + Cross-modal
};
\end{tikzpicture}
\caption{Espacio metodológico de métodos de \UQ.
Ningún método existente combina las tres propiedades simultáneamente.
Nuestra contribución (KL Cross-Modal) ocupa el octante vacío.}
\label{fig:espacio-metodologico}
\end{figure}

Como se ilustra en la \Cref{fig:espacio-metodologico}, \textbf{ningún método publicado combina las tres propiedades} de ser single-pass, basado en features internas y cross-modal.
Esta es la novedad metodológica central de nuestra propuesta.

% ============================================================
\section{Método Propio}
% ============================================================

\subsection{Intuición}

Cuando un modelo de visión-lenguaje ``ve'' algo en la imagen pero ``dice'' otra cosa en su respuesta, ese desacuerdo interno se puede medir con un número.
Ese número nos dice cuándo el modelo está a punto de equivocarse, sin necesidad de saber la respuesta correcta.

La analogía es la del cerebro humano: si los ojos ven a una persona decir ``ga'' pero los oídos oyen ``ba'', el cerebro detecta el \textbf{conflicto entre sentidos} y se pone en alerta.
Nosotros construimos el equivalente computacional: medimos el desacuerdo entre la representación interna de la imagen y la representación interna de la respuesta del modelo, con la divergencia de Kullback-Leibler (\KL).

\subsection{Arquitectura del modelo base}

Utilizamos \texttt{google/medgemma-4b-it} \citep{sellergren2025medgemma}, un modelo de visión-lenguaje médico con la siguiente arquitectura:

\begin{table}[H]
\centering
\caption{Especificaciones de MedGemma 4B}
\label{tab:medgemma-specs}
\begin{tabular}{ll}
\toprule
\textbf{Componente} & \textbf{Especificación} \\
\midrule
Parámetros totales & $\sim$4B \\
Vision encoder & MedSigLIP (SigLIP-400M fine-tuned médico) \\
Capas vision encoder & 27 \\
Hidden size vision & 1152 \\
Decoder & Gemma 3, 34 capas transformer \\
Hidden size decoder & 2560 \\
Vocabulario & 262,144 tokens \\
Input resolution & 896$\times$896 \\
Patch size & 14$\times$14 \\
Tokens visuales & 256 (tras pooling 4$\times$4) \\
VRAM (bfloat16) & $\sim$10--12 GB \\
VRAM (4-bit NF4) & $\sim$4--5 GB \\
\bottomrule
\end{tabular}
\end{table}

\subsection{Extracción de representaciones}

\subsubsection{Tokens visuales}

La imagen se procesa mediante el vision encoder MedSigLIP:

\begin{enumerate}
    \item Resize a 896$\times$896 píxeles (realizado por el procesador del modelo).
    \item División en parches de 14$\times$14: $896/14 = 64$ parches por lado $\rightarrow$ $64 \times 64 = 4{,}096$ parches.
    \item Cada parche pasa por la torre de visión (SigLIP) y sale como vector de 1,152 dimensiones.
    \item Pooling 4$\times$4 agrupa parches vecinos: $4{,}096 / 16 = 256$ vectores.
    \item Proyección a 2,560 dimensiones (hidden size del decoder).
\end{enumerate}

Los 256 tokens visuales se insertan en la secuencia del decoder como si fueran 256 palabras adicionales, identificables por el token especial \texttt{<image\_soft\_token>} (ID 262,144).

\subsubsection{Hidden states}

El decoder es una pila de 34 capas transformer idénticas.
Al pedir \texttt{output\_hidden\_states=True}, \texttt{generate()} devuelve una tupla con 35 entradas por paso de generación:

\begin{itemize}
    \item Índice 0: embeddings de entrada (no usados en la KL).
    \item Índices 1--34: salidas de las 34 capas.
\end{itemize}

Con \texttt{max\_new\_tokens=1}, solo existe el prefill (\texttt{hidden\_states[0]}).
El estado que condiciona la respuesta es la \textbf{última posición del prefill}: \texttt{hidden\_states[0][capa][:, -1, :]}.

\subsubsection{Tres ingredientes por imagen}

\begin{table}[H]
\centering
\caption{Ingredientes extraídos en una sola pasada}
\label{tab:ingredientes}
\begin{tabularx}{\textwidth}{lXl}
\toprule
\textbf{Variable} & \textbf{Descripción} & \textbf{Origen exacto} \\
\midrule
$p_{\text{vis}}$ & Resumen de ``lo que el modelo vio'' & 256 hidden states de tokens de imagen (máscara por ID 262,144), mean-pooling \\
$p_{\text{text}}$ & ``Lo que el modelo está a punto de decir'' & Hidden state de la última posición del prefill \\
$p_{\text{yes}}$, $p_{\text{no}}$ & La respuesta concreta & Logits de ``yes'' (ID 4443) y ``no'' (ID 1904) en \texttt{scores[0]} \\
\bottomrule
\end{tabularx}
\end{table}

\subsection{De vector a distribución: softmax temperado}

La \KL\ compara distribuciones de probabilidad.
Convertimos los vectores de hidden states en distribuciones mediante softmax temperado:

\begin{equation}
p_i = \frac{\exp(z_i / \tau)}{\sum_j \exp(z_j / \tau)}
\end{equation}

donde $\tau$ es la temperatura.
Los hidden states de Gemma tienen componentes de magnitud $10^2$--$10^3$ (\textit{massive activations}).
Con $\tau = 1$, el softmax en float32 colapsa 2,559 de 2,560 componentes a exactamente 0, produciendo $\KL = \infty$.
Por ello, computamos \texttt{log\_softmax} en \textbf{float64}, lo que garantiza valores finitos y $\KL(p \| p) = 0$ exacto.

\subsection{Divergencia KL}

La divergencia de Kullback-Leibler se define como:

\begin{equation}
\KL(p \| q) = \sum_i p_i \cdot \ln\left(\frac{p_i}{q_i}\right)
\end{equation}

\textbf{Propiedades relevantes:}

\begin{itemize}
    \item $\KL(p \| q) \geq 0$, con igualdad solo si $p = q$.
    \item \textbf{Asimétrica:} $\KL(p \| q) \neq \KL(q \| p)$.
    Esta asimetría es una \textit{feature}: cada dirección pregunta algo distinto.
    \item No está acotada superiormente (puede valer hasta infinito).
\end{itemize}

Computamos tres variantes:

\begin{align}
u_{\text{vis}\rightarrow\text{text}}(x) &= \KL(p_{\text{vis}} \| p_{\text{text}}) \\
u_{\text{text}\rightarrow\text{vis}}(x) &= \KL(p_{\text{text}} \| p_{\text{vis}}) \\
u_{\text{JSD}}(x) &= \frac{1}{2}\KL\left(p_{\text{vis}} \Big\| \frac{p_{\text{vis}} + p_{\text{text}}}{2}\right) + \frac{1}{2}\KL\left(p_{\text{text}} \Big\| \frac{p_{\text{vis}} + p_{\text{text}}}{2}\right)
\end{align}

\subsection{Interpretación direccional}

La asimetría de la \KL\ proporciona diagnóstico estructural del modo de fallo:

\begin{table}[H]
\centering
\caption{Interpretación de las direcciones de \KL}
\label{tab:interpretacion-kl}
\begin{tabularx}{\textwidth}{XX}
\toprule
\textbf{Señal} & \textbf{Interpretación clínica} \\
\midrule
$\KL(\text{vis} \| \text{text})$ alta & El texto ignora la imagen: \textit{language prior bias}, alucinación textual \\
$\KL(\text{text} \| \text{vis})$ alta & La imagen contiene hallazgos atípicos no capturados por el vocabulario \\
Ambas altas & Ambigüedad clínica intrínseca (caso genuinamente difícil) \\
Ambas bajas & Acuerdo confidente (caso típico, bien representado) \\
\bottomrule
\end{tabularx}
\end{table}

\subsection{Algoritmo completo}

\begin{algorithm}[H]
\caption{Pipeline de extracción de señal de incertidumbre}
\label{alg:pipeline}
\begin{algorithmic}[1]
\Require Imagen de fondo de ojo $x$, prompt $p$, modelo MedGemma $M$
\Ensure Señal de incertidumbre $u(x)$, predicción $\hat{y}$, probabilidad $P(\text{yes})$
\State Procesar $x$ a 896$\times$896 vía \texttt{AutoProcessor}
\State Armar secuencia con chat template: imagen + prompt $p$
\State $\text{outputs} \gets M.\text{generate}(\text{max\_new\_tokens}=1, \text{output\_hidden\_states}=\text{True}, \text{output\_scores}=\text{True})$
\State Extraer \texttt{scores[0]}: logit\_yes, logit\_no
\State $P(\text{yes}) \gets \frac{\exp(\text{logit\_yes})}{\exp(\text{logit\_yes}) + \exp(\text{logit\_no})}$
\State $\hat{y} \gets \mathbb{1}[P(\text{yes}) > 0.5]$
\For{capa $\ell \in \{17, 26, 34\}$}
    \State $h_{\text{prefill}} \gets \text{outputs}.\text{hidden\_states}[0][\ell]$  \Comment{35 entradas, índice 0 = embeddings}
    \State $\text{mask} \gets (\text{input\_ids} == 262144)$  \Comment{256 tokens de imagen}
    \State $p_{\text{vis}} \gets \text{mean\_pool}(h_{\text{prefill}}[\text{mask}])$
    \State $p_{\text{text}} \gets h_{\text{prefill}}[:, -1, :]$  \Comment{última posición del prefill}
    \For{$\tau \in \{1, 2, 4\}$}
        \State $P_{\text{vis}} \gets \text{softmax}(p_{\text{vis}} / \tau)$  \Comment{en float64}
        \State $P_{\text{text}} \gets \text{softmax}(p_{\text{text}} / \tau)$  \Comment{en float64}
        \State Guardar $\KL(P_{\text{vis}} \| P_{\text{text}})$, $\KL(P_{\text{text}} \| P_{\text{vis}})$, $\JSD(P_{\text{vis}}, P_{\text{text}})$
    \EndFor
\EndFor
\State Guardar fila CSV con 54 variantes + logits + predicción
\end{algorithmic}
\end{algorithm}

% ============================================================
\section{Diseño Experimental}
% ============================================================

\subsection{Dataset}

Utilizamos MM-ODIR-129 \citep{mmodir129}, un dataset de 129 fotografías de fondo de ojo completas (resolución variable), originarias de ODIR-5K y re-anotadas por oftalmólogos de Costa Rica.

\begin{table}[H]
\centering
\caption{Composición del dataset MM-ODIR-129}
\label{tab:dataset}
\begin{tabular}{lccc}
\toprule
\textbf{Split} & \textbf{Normal} & \textbf{Patológico (Glaucoma)} & \textbf{Total} \\
\midrule
Train & 36 & 41 & 77 \\
Validation & 12 & 14 & 26 \\
Test & 12 & 14 & 26 \\
\midrule
\textbf{Total} & \textbf{60} & \textbf{69} & \textbf{129} \\
\bottomrule
\end{tabular}
\end{table}

Cada imagen patológica incluye 7 gradings ordinales de glaucoma, incluyendo \texttt{cup\_to\_disc\_ratio} (ordinal 0--4), que utilizamos para la hipótesis H4.
Las imágenes patológicas cuentan con máscaras de copa y disco óptico (PNG + NPY), reservadas para trabajo futuro.

\subsection{Prompts}

\begin{itemize}[leftmargin=2em]
    \item \textbf{P1 (principal):} ``Does this fundus image show glaucoma? Answer yes or no.''
    
    \item \textbf{P4 (contraste):} system prompt ``You are an expert ophthalmologist.'' + la misma pregunta.
\end{itemize}

\subsection{Baselines}

Comparamos nuestra señal contra los siguientes métodos:

\begin{table}[H]
\centering
\caption{Baselines de comparación}
\label{tab:baselines}
\begin{tabularx}{\textwidth}{lXl}
\toprule
\textbf{Método} & \textbf{Descripción} & \textbf{Costo} \\
\midrule
MSP & Maximum Softmax Probability sobre la respuesta & Single-pass \\
Entropía & Entropía de la distribución yes/no & Single-pass \\
TS + Entropía & Temperature Scaling ajustado en train, luego entropía & Single-pass \\
Energy Score & Score de energía de \citet{liu2020energy} & Single-pass \\
Mahalanobis & Distancia de Mahalanobis en espacio de features & Single-pass \\
Self-Consistency & 10 muestras a $T=0.7$, entropía de la fracción de ``yes'' & Multi-pass (10$\times$) \\
\bottomrule
\end{tabularx}
\end{table}

\subsection{Métricas}

\begin{itemize}[leftmargin=2em]
    \item \textbf{\AUROC} (primaria): Probabilidad de que un error tenga mayor $u(x)$ que un acierto.
    
    \item \textbf{\AUPRC} (secundaria): Precision-recall curve para detección de errores.
    
    \item \textbf{Brier Score}: Calibración probabilística de $P(\text{yes})$.
    
    \item \textbf{\ECE}: Expected Calibration Error.
    
    \item \textbf{Sensitivity@80\%Specificity}: Umbral clínico operativo.
    
    \item \textbf{Accuracy vs.~Coverage}: Curva de clasificación selectiva.
\end{itemize}

\subsection{Análisis estadístico}

\begin{itemize}[leftmargin=2em]
    \item \textbf{Mann-Whitney U}: ¿Los errores tienen sistemáticamente mayor $u(x)$ que los aciertos?
    Reportamos effect size $r = |Z| / \sqrt{N}$ (convención de Fritz et al., 2012).
    
    \item \textbf{Bootstrap BCa}: Intervalos de confianza del 95\% para \AUROC\ y \AUPRC\ con 9,999 remuestreos.
    BCa corrige sesgo y asimetría, crítico con $N=129$.
    
    \item \textbf{DeLong's test} (exploratorio): Comparación significativa de curvas \AUROC\ contra el mejor baseline.
    
    \item \textbf{Spearman + permutación} (H4): Correlación entre $u(x)$ y \texttt{cup\_to\_disc\_ratio} en los 69 patológicos.
    El p-value por permutación (9,999 barajadas) es exacto sin asumir normalidad.
\end{itemize}

\subsection{Variantes y ablaciones}

Por imagen computamos \textbf{54 variantes} en una sola pasada:

\begin{equation}
3 \text{ divergencias} \times 3 \text{ capas} \times 3 \text{ temperaturas} \times 2 \text{ poolings} = 54 \text{ variantes}
\end{equation}

La variante reportada se selecciona \textbf{exclusivamente en el split de entrenamiento} (77 imágenes) para evitar contaminación estadística.
Todo lo reportado en test usa esa variante congelada.

\subsection{Reproducibilidad}

\begin{itemize}[leftmargin=2em]
    \item Seeds fijadas: 42 para random, numpy, torch.
    \item \texttt{torch.backends.cudnn.deterministic = True}.
    \item \texttt{CUBLAS\_WORKSPACE\_CONFIG=:4096:8}.
    \item Greedy decoding (determinista) para la pasada principal.
    \item Transformers $\geq$ 4.51.3 (obligatorio para API de hidden states).
    \item Repositorio anónimo en \texttt{anonymous.4open.science} para double-blind review.
\end{itemize}

% ============================================================
\section{Resultados Esperados y Discusión}
% ============================================================

\subsection{Criterios de éxito}

Definimos tres escenarios posibles según el \AUROC\ de detección de errores:

\begin{table}[H]
\centering
\caption{Escenarios de resultado y su interpretación}
\label{tab:escenarios}
\begin{tabularx}{\textwidth}{lXX}
\toprule
\textbf{Escenario} & \textbf{Lectura} & \textbf{Consecuencia} \\
\midrule
\AUROC $\geq$ 0.70, significativo & Éxito pleno & Paper fuerte: señal útil, barata y novedosa \\
\AUROC 0.60--0.70 & Éxito parcial & Publicable en BIP: comparable a baselines baratos a costo mínimo \\
\AUROC $\approx$ 0.5 & Resultado negativo & Plan B: análisis de \textit{por qué} no funciona + estudio de limitaciones \\
\bottomrule
\end{tabularx}
\end{table}

\subsection{Costo computacional}

\begin{table}[H]
\centering
\caption{Comparación de costos computacionales por imagen (GPU 16GB)}
\label{tab:costos}
\begin{tabular}{lcccc}
\toprule
\textbf{Método} & \textbf{Passes} & \textbf{VRAM} & \textbf{Tiempo} & \textbf{Training} \\
\midrule
MC-Dropout (25) & 25 & 16 GB & $\sim$60 s & No \\
Semantic Entropy (10) & 10+NLI & 16 GB+ & $\sim$30 s & No \\
Deep Ensemble (5) & 5 & 80 GB & $\sim$12 s & Sí (5$\times$) \\
SAPLMA Probe & 1 & 16 GB & $\sim$2 s & Sí (probe) \\
\midrule
\textbf{KL Cross-Modal} & \textbf{1} & \textbf{4--5 GB} & \textbf{$\sim$2 s} & \textbf{No} \\
\bottomrule
\end{tabular}
\end{table}

El costo computacional mínimo hace el método viable para deployment en el Global South (hospitales rurales, Latinoamérica), donde los recursos son limitados.

\subsection{Implicaciones clínicas}

La señal KL cross-modal habilita un sistema de \textbf{triage automatizado}:

\begin{enumerate}
    \item El modelo diagnostica glaucoma y emite $P(\text{yes})$.
    \item Simultáneamente, computa $u(x) = \KL(p_{\text{vis}} \| p_{\text{text}})$.
    \item Si $u(x)$ excede un umbral, el caso se deriva automáticamente a revisión humana.
    \item Esto enriquece el sistema \texttt{ophthalmo\_capture} con una capa de confiabilidad.
\end{enumerate}

\subsection{Conexión con los pilares de la tesis}

\subsubsection{UQ $\rightarrow$ XAI: La incertidumbre como explicación nativa}

La magnitud y dirección de la \KL\ son inherentemente explicativas.
A diferencia de Grad-CAM o SHAP (métodos \textit{post-hoc} que requieren cálculos adicionales), la \KL\ es una medida \textit{native} de la arquitectura VLM.
No se explica \textit{qué} predice el modelo, sino \textit{por qué} el modelo no está seguro.

\subsubsection{UQ $\rightarrow$ Few-Shot: La incertidumbre como curador de datos}

En few-shot, 10 ejemplos con alta \KL\ son peores que 5 ejemplos con baja \KL.
La \KL\ puede funcionar como \textbf{curador de support sets}: seleccionar ejemplos que maximicen la coherencia cross-modal.

\subsubsection{UQ $\rightarrow$ Segmentación: La incertidumbre como mapa de riesgo espacial}

La \KL\ global puede descomponerse en contribuciones locales mediante técnicas de \textit{attention attribution}, generando un \textbf{mapa de riesgo espacial} donde las regiones de alta contribución a la \KL\ son precisamente las regiones anatómicas ambiguas.

% ============================================================
\section{Conclusiones}
% ============================================================

\subsection{Síntesis}

Este informe presentó la propuesta de investigación para BIP 2026: una señal de incertidumbre \textit{lightweight} basada en la divergencia \KL\ cross-modal en MedGemma 4B para detección de glaucoma.
Las contribuciones principales son:

\begin{enumerate}
    \item \textbf{Novedad metodológica:} Primer método que combina single-pass + feature-based + cross-modal para \UQ\ en \MedVLM s oftalmológicos.
    
    \item \textbf{Justificación bio-inspirada:} Formalización computacional del principio neurocientífico de detección de conflictos cross-modal.
    
    \item \textbf{Interpretabilidad direccional:} La asimetría de \KL\ diagnostica el \textit{modo de fallo} (language prior bias vs.~ambigüedad visual vs.~ambigüedad clínica).
    
    \item \textbf{Escalabilidad:} 1 pasada, 4--5 GB VRAM, $\sim$2 segundos por imagen: viable para despliegue en recursos limitados.
\end{enumerate}

\subsection{Trabajo futuro}

\begin{enumerate}
    \item \textbf{Calibración:} Transformar \KL\ en probabilidad predictiva calibrada (\ECE).
    
    \item \textbf{Triage adaptativo:} Activar MC-Dropout solo cuando $\KL$ es alta, reduciendo costo promedio.
    
    \item \textbf{Segmentación text-guided:} Extender la señal \KL\ a segmentación OD/OC con prompts descriptivos.
    
    \item \textbf{Few-shot con filtrado KL:} Desarrollar algoritmo de selección de support sets basado en coherencia cross-modal.
    
    \item \textbf{Validación clínica prospectiva:} Evaluación en cohorte real con métricas de human-AI interaction.
    
    \item \textbf{Generalización cross-dataset:} Validar en REFUGE, RIM-ONE y PAPILA.
\end{enumerate}

% ============================================================
\section*{Agradecimientos}
% ============================================================

Este trabajo se enmarca dentro del proyecto de doctorado en Inteligencia Artificial Aplicada a Oftalmología del Instituto Tecnológico de Costa Rica, bajo la tutoría del Dr.~Saúl Calderón Ramírez.
Agradecemos a los oftalmólogos que contribuyeron con las anotaciones del dataset MM-ODIR-129.

% ============================================================
% REFERENCIAS
% ============================================================
\newpage
\bibliographystyle{apalike}
\begin{thebibliography}{99}

\bibitem[Azaria y Mitchell, 2023]{azaria2023internal}
Azaria, A. y Mitchell, T. (2023).
The internal state of an LLM knows when its lying.
\emph{Findings of EMNLP}, 967--976.

\bibitem[Badash et~al., 2026]{badash2026between}
Badash, Z.~N., Belinkov, Y., y Freiman, M. (2026).
Between the layers lies the truth: Compact, per-instance uncertainty via cross-layer agreement.
\emph{arXiv preprint arXiv:2603.22299}.

\bibitem[Botvinick et~al., 2001]{botvinick2001conflict}
Botvinick, M.~M., Braver, T.~S., Barch, D.~M., Carter, C.~S., y Cohen, J.~D. (2001).
Conflict monitoring and cognitive control.
\emph{Psychological Review}, 108(3), 624--652.

\bibitem[Ernst y B\"ulthoff, 2004]{ernst2004integration}
Ernst, M.~O. y B\"ulthoff, H.~H. (2004).
Merging the senses into a robust percept.
\emph{Trends in Cognitive Sciences}, 8(4), 162--169.

\bibitem[Gal y Ghahramani, 2016]{gal2016dropout}
Gal, Y. y Ghahramani, Z. (2016).
Dropout as a bayesian approximation: Representing model uncertainty in deep learning.
\emph{ICML}, 1050--1059.

\bibitem[Hendrycks y Gimpel, 2017]{hendrycks2017baseline}
Hendrycks, D. y Gimpel, K. (2017).
A baseline for detecting misclassified and out-of-distribution examples in neural networks.
\emph{ICLR}.

\bibitem[Kuhn et~al., 2023]{kuhn2023semantic}
Kuhn, L., Gal, Y., y Farquhar, S. (2023).
Semantic uncertainty: Linguistic invariances for uncertainty estimation in natural language generation.
\emph{ICLR}.

\bibitem[Lakshminarayanan et~al., 2017]{lakshminarayanan2017ensembles}
Lakshminarayanan, B., Pritzel, A., y Blundell, C. (2017).
Simple and scalable predictive uncertainty estimation using deep ensembles.
\emph{NeurIPS}, 30.

\bibitem[Lee et~al., 2018]{lee2018mahalanobis}
Lee, K., Lee, H., Lee, K., y Shin, J. (2018).
Training confidence-calibrated classifiers for detecting out-of-distribution samples.
\emph{ICLR}.

\bibitem[Liang et~al., 2025]{liang2025uncertainty}
Liang, X., Zhang, Y., et~al. (2025).
Uncertainty-driven expert control: Enhancing the reliability of medical vision-language models.
\emph{ICCV}.

\bibitem[Liu et~al., 2020]{liu2020energy}
Liu, W., Wang, X., Owens, J., y Li, Y. (2020).
Energy-based out-of-distribution detection.
\emph{NeurIPS}, 33, 21464--21475.

\bibitem[McDermott et~al., 2024]{mcdermott2024auroc}
McDermott, J., Wang, S., et~al. (2024).
A closer look at \AUROC\ and \AUPRC\ under class imbalance.
\emph{NeurIPS}, 37, 44102--44163.

\bibitem[McGurk y MacDonald, 1976]{mcgurk1976}
McGurk, H. y MacDonald, J. (1976).
Hearing lips and seeing voices.
\emph{Nature}, 264, 746--748.

\bibitem[Saito y Rehmsmeier, 2015]{saito2015precision}
Saito, T. y Rehmsmeier, M. (2015).
The precision-recall plot is more informative than the ROC plot when evaluating binary classifiers on imbalanced datasets.
\emph{PLOS ONE}.

\bibitem[Sellergren et~al., 2025]{sellergren2025medgemma}
Sellergren, A., Azizi, S., et~al. (2025).
MedGemma: Technical report.
\emph{arXiv preprint arXiv:2507.05201}.

\bibitem[Shi et~al., 2025]{shietal2025eyeclip}
Shi, D., et~al. (2025).
EyeCLIP: A multimodal visual-language foundation model for computational ophthalmology.
\emph{Nature Medicine / arXiv:2409.06644}.

\bibitem[Yeung et~al., 2004]{yeung2004conflict}
Yeung, N., Botvinick, M.~M., y Cohen, J.~D. (2004).
The neural basis of error detection: Conflict monitoring and the error-related negativity.
\emph{Psychological Review}, 111(4), 931--959.

\bibitem[Modality Bias, 2025]{modalitybias2025}
(2025).
Clinical decision-making relies on the integrated analysis of medical images and associated textual information.
\emph{arXiv:2508.00171}.

\bibitem[VLM-UQBench, 2026]{vlmuqbench2026}
(2026).
VLM-UQBench: A benchmark for uncertainty quantification in vision-language models.
\emph{arXiv:2602.09214}.

\bibitem[UMPIRE, 2026]{umpire2026}
(2026).
UMPIRE: Uncertainty estimation via multimodal perturbation and inconsistency-aware representation.
\emph{arXiv:2602.24195}.

\bibitem[VIG-TUQ, 2026]{vigtuq2026}
(2026).
Visual-grounded token uncertainty quantification for vision-language models.
\emph{arXiv:2605.27136}.

\bibitem[MM-ODIR-129, 2026]{mmodir129}
(2026).
MM-ODIR-129: A curated dataset for glaucoma detection in fundus images.
\emph{HuggingFace Datasets}.

\bibitem[VOLMO, 2026]{volmo2026}
(2026).
VOLMO: Versatile and open large models for ophthalmology.
\emph{arXiv:2603.23953}.

\bibitem[RetiZero, 2024]{retizero2024}
(2024).
RetiZero: Common and rare fundus diseases identification using vision-language foundation model.
\emph{arXiv:2406.09317}.

\end{thebibliography}

\end{document}
